% =====================================================================
%  Registro do Processo de Criação — Stage 6 (academic-pipeline)
%  Colaboração humano-IA · "Estrutura Produtiva do Espírito Santo"
% =====================================================================
\documentclass[12pt,a4paper]{article}
\usepackage[utf8]{inputenc}
\usepackage[T1]{fontenc}
\usepackage[brazil]{babel}
\usepackage{booktabs}
\usepackage{longtable}
\usepackage{geometry}
\usepackage{xcolor}
\usepackage{fancyhdr}
\usepackage[hidelinks]{hyperref}
\geometry{margin=2.5cm}
\definecolor{barcol}{RGB}{120,120,120}
\newcommand{\scorebar}[2]{{\color{barcol}\rule{#1mm}{2.6mm}}\;\textbf{#2}}

\pagestyle{fancy}\fancyhf{}
\lhead{\itshape Registro do Processo de Criação}
\rhead{26 jun.\ 2026}
\cfoot{\thepage}
\renewcommand{\headrulewidth}{0.4pt}

\title{\textbf{Registro do Processo de Criação}\\[4pt]
\large ``Estrutura Produtiva do Espírito Santo: uma análise de insumo-produto''\\[2pt]
\normalsize Stage 6 do pipeline \texttt{academic-pipeline} · colaboração humano-IA (Claude Code / Opus 4.8)}
\author{Felipe Carvalho — PPGEco/UFES \quad$\cdot$\quad registro gerado pela IA}
\date{26 de junho de 2026}

\begin{document}
\maketitle
\thispagestyle{fancy}
\tableofcontents
\bigskip\hrule\bigskip

% =====================================================================
\section{Informações do artigo}
\begin{longtable}{p{4.3cm}p{10cm}}
\toprule
\textbf{Título final} & Estrutura Produtiva do Espírito Santo: uma análise de insumo-produto \\
\textbf{Subtítulo} & Multiplicadores, setores-chave, vocação territorial e abertura de uma economia de base (2008--2021) \\
\textbf{Autor} & Felipe Carvalho — PPGEco/UFES \\
\textbf{Entregável} & \texttt{paper/es\_estrutura\_produtiva.pdf} (10 págs · 4 figuras · 2 tabelas · 17 refs) \\
\textbf{Pipeline} & \texttt{pesquisa/14--22} + \texttt{diag\_sider.py}; CSVs em \texttt{outputs/}; figuras em \texttt{figuras/} \\
\textbf{Irmão arquivado} & \texttt{es\_plataforma\_fractal.tex} (tese de invariância de escala, depois absorvida) \\
\textbf{Git} & commit \texttt{1ed4101} em \texttt{origin/main} (github.com/fcarva/es-insumo-produto) \\
\bottomrule
\end{longtable}

\noindent\textbf{Instrução inicial (verbatim):} \emph{``fazer um grande review do projeto em fase das
novas questões\dots\ ter um olhar panorâmico mais amplo\dots\ fizesse uma deep research com base no que
já tem\dots\ pegasse as matrizes novas e ver o que pode ser feito use full mode da skill de ponta a
ponta.''}

% =====================================================================
\section{Processo estágio a estágio}
O pipeline rodou \textbf{duas vezes}: primeiro sobre uma tese de invariância de escala (descartada pelo
autor), depois sobre a caracterização (entregue).

\subsection{Ciclo A — tese ``economia-plataforma'' (abandonado a pedido)}
\begin{longtable}{p{2.6cm}p{8.3cm}p{3.4cm}}
\toprule
\textbf{Estágio} & \textbf{Entrada $\to$ saída} & \textbf{Decisão do autor} \\
\midrule
1 RESEARCH & acervo + matrizes $\to$ \texttt{DEEP\_RESEARCH\_PANORAMA} (6 RQs multiescala) & RQ-A + RQ-B (fractal + benchmark) \\
2 análise & \texttt{14/15/16} $\to$ achado \emph{honesto}: mecânica genérica/invariante, especificidade composicional (contrariou a hipótese) & reenquadrar + endurecer antes \\
2 WRITE & $\to$ \texttt{es\_plataforma\_fractal.tex} & — \\
2.5 INTEGRITY & \textbf{PASS}, mas o gate pegou erro factual (``Mineração dominante em 2 estados'' $\to$ \textbf{3}: PA, ES, RJ); corrigido & pausa para ler \\
\bottomrule
\end{longtable}

\subsection{Ciclo B — caracterização (entregue)}
\noindent\textbf{Correção de rumo (verbatim):} \emph{``Eu não gostei muito do título economia fractal\dots\
fizesse uma outra deep search, pra gente fazer uma caracterização do espírito santo\dots\ estudasse a
partir do Nereus\dots\ como o Haddad\dots\ Esse novo artigo, ele ficou interessante mas não é essa a
pegada que eu quero fazer. Vamos voltar.''}

\begin{longtable}{p{2.6cm}p{8.3cm}p{3.4cm}}
\toprule
\textbf{Estágio} & \textbf{Entrada $\to$ saída} & \textbf{Decisão do autor} \\
\midrule
1 RESEARCH & papers MT/RS + lit.\ nacional $\to$ \texttt{DEEP\_RESEARCH\_CARACTERIZACAO} (gênero estadual; dados NEREUS) & backbone C1+C2 + título do gênero \\
2 (C1) & \texttt{17} $\to$ retrato 2008 (mult.\ prod/emprego/renda I e II, RH, ligações puras, setores-chave) & — \\
2 (C2) & \texttt{18} $\to$ trajetória nacional 2010--2021 (Nível 68) & adicionar C3 antes \\
2 (C3) & \texttt{20/21} $\to$ vocação das 10 microrregiões (quociente locacional) & — \\
2 WRITE & $\to$ \texttt{es\_estrutura\_produtiva.tex} & — \\
2.5 INTEGRITY & \textbf{PASS} (0 erros) & checar siderurgia 2021 \\
robustez & \texttt{diag\_sider.py} $\to$ não é quebra de classificação, é \textbf{efeito de preço} (aço 2021, preços correntes) & Stage 3 REVIEW \\
3 REVIEW & parecer 5-revisores $\to$ \textbf{Major Revision} (4 obrigatórios) & atacar os 4 itens \\
4 REVISE & R1 benchmark (\texttt{22}) · R2 objeto temporal · R3 lit.\ ES (Ribeiro 2024; Sessa 2017) · R4 política & re-review + finalizar \\
3' RE-REVIEW & Minor $\to$ Accept; pegou \textbf{NEW-1} (1,76 vs 1,66) $\to$ corrigido & — \\
4.5 FINAL INTEGRITY & \textbf{PASS (zero issues)} & — \\
5 FINALIZE & PDF compilado do LaTeX & — \\
\bottomrule
\end{longtable}

\subsection{Pós-finalização (refinamentos do autor)}
\begin{itemize}\setlength\itemsep{1pt}
\item \textbf{Sugeridos S1--S4} atacados (correlação PTL$\times$VBP; tipo II; \emph{doença holandesa};
  padrão não-óbvio do C3 = enclave territorial).
\item \textbf{Deep research} para ancorar o S3 na literatura \emph{nacional} de IO (\emph{``o melhor do
  insumo''}): Morceiro (2012); Nassif et al.\ (2015, citado honestamente como cético); Oreiro \& Feijó (2010).
\item \textbf{Conciliação} com a ``economia-plataforma'': incorporada como §7 (abertura/spillover-feedback
  + escalas aninhadas), reconciliando base (estrutura) + plataforma (abertura).
\item \textbf{Figuras}: fundo transparente + boas práticas (skill \texttt{data:data-visualization}).
\item \textbf{Terminologia}: ``fractal'' $\to$ \textbf{``núcleo-periferia aninhada / escalas aninhadas''}
  (Friedmann, 1966).
\end{itemize}

% =====================================================================
\section{Resumo do padrão de interação}
\begin{longtable}{p{8.5cm}p{6cm}}
\toprule
\textbf{Métrica} & \textbf{Valor} \\
\midrule
Ciclos de pipeline & 2 (tese descartada + caracterização) \\
Deep researches & 3 (panorama · caracterização · S3) \\
Checkpoints com decisão do autor & 10 (todos respondidos) \\
Correção de rumo maior & \textbf{1} (o pivô) \\
Rodadas de parecer & 1 full (5 papéis) + 1 re-review \\
Verificações de integridade & 4 \\
Erros pegos pelos gates & 2 (``dois únicos''$\to$três; NEW-1) + 1 anomalia investigada (2021) \\
Scripts criados & 10 (\texttt{14--22}, \texttt{diag\_sider}) \\
Referências finais (verificadas) & 17 \\
Papel do autor & diretor editorial e curador intelectual \\
Papel da IA & pesquisa, código, análise, redação, parecer, verificação \\
\bottomrule
\end{longtable}

% =====================================================================
\section{Decisões-chave do autor (cronológicas)}
\begin{enumerate}\setlength\itemsep{1pt}
\item Rodar o pipeline em \textbf{full mode}, panorâmico, sobre as matrizes novas.
\item Desenvolver \textbf{RQ-A + RQ-B} (fractal + benchmark).
\item \textbf{Reenquadrar a tese} + \textbf{endurecer a análise} antes de redigir.
\item \textbf{PIVÔ:} abandonar a pegada ``fractal''; fazer uma \textbf{caracterização} do ES via IO,
  ancorada na escola brasileira (NEREUS, Haddad). \emph{(a decisão mais importante da sessão.)}
\item \textbf{Backbone C1 + C2} + título no gênero clássico.
\item \textbf{Adicionar C3} (microrregiões) antes de redigir.
\item \textbf{Investigar} a queda da siderurgia em 2021 antes de afirmá-la.
\item Submeter ao \textbf{parecer 5-revisores}.
\item \textbf{Atacar os 4 itens obrigatórios}.
\item \textbf{Re-review + finalizar}.
\item \textbf{Atacar os sugeridos S1--S4}.
\item \textbf{Deep research} para ancorar o S3 na literatura nacional de IO.
\item \textbf{Conciliar} a economia-plataforma como face complementar.
\item \textbf{Figuras} sem fundo, em boas práticas.
\item \textbf{Trocar ``fractal''} por termo da literatura; \textbf{commit + push}; iniciar \textbf{Stage 6}.
\end{enumerate}

% =====================================================================
\section{Lições-chave (reutilizáveis)}
\begin{enumerate}\setlength\itemsep{1pt}
\item \textbf{Um pivô precoce vale mais que um polimento tardio.} Trocar a pegada custou um ciclo, mas
  levou a um artigo mais defensável e alinhado à tradição do campo.
\item \textbf{O benchmark salva a tese da tautologia.} Só o contraste com as 27 UFs transformou o ES em
  ``caso-limite'' — afirmação identificada, não retórica.
\item \textbf{Anomalia de dado merece investigação, não nota apressada.} A queda da siderurgia-2021 era
  efeito de preço (matrizes a preços correntes), não quebra estrutural.
\item \textbf{Citação só depois de verificada.} As 17 referências foram conferidas antes de entrar.
\item \textbf{Conciliar $>$ descartar.} O trabalho ``fractal'' abandonado virou a §7, enriquecendo o artigo.
\end{enumerate}

% =====================================================================
\section{Relatório de autorreflexão da IA}
\noindent\emph{Ironia registrada (princípio do protocolo): esta autorreflexão é produzida pela mesma IA
que pode ter sido condescendente durante o pipeline. Leia com essa consciência.}

\subsection*{Resumo comportamental}
A IA operou como executora-sob-checkpoint: pesquisou, codou, analisou, redigiu, criticou (parecer
adversarial real) e verificou. Em dois momentos reportou achados que \textbf{contrariaram a própria
hipótese} (o benchmark ``genérico, não especial''; a queda 2021 = preço) — evidência contra
condescendência substantiva.

\subsection*{Risco de sicofância: BAIXO}
0 alertas de saúde; o parecer foi genuinamente exigente (Major Revision); achados desconfortáveis foram
reportados sem suavização.

\subsection*{Frame-lock — 1 incidente real}
A IA enquadrou o trabalho inicial como \textbf{teste de tese} (``fractal'') e \textbf{não detectou
sozinha} que o autor queria uma \textbf{caracterização}. Foi preciso o pivô do autor para corrigir. É o
limite mais honesto da sessão: a IA otimizou dentro do frame que ela mesma propôs.

\subsection*{O que a IA errou (gates funcionando)}
\begin{itemize}\setlength\itemsep{1pt}
\item \textbf{``Dois únicos'' $\to$ três} (Mineração dominante): erro factual pego no Stage 2.5.
\item \textbf{NEW-1}: a própria revisão introduziu inconsistência numérica (1,76 simples vs 1,66
  ponderado), pega no re-review e reconciliada.
\item \textbf{Frame inicial}: precisou de correção humana.
\item \textbf{Hipótese inicial da RQ-B} (ES ``especial''): refutada pelo dado, reportada como tal.
\end{itemize}

\subsection*{Log de modos de falha (Stage 2.5/4.5)}
\begin{longtable}{p{6.2cm}p{8cm}}
\toprule
\textbf{Modo} & \textbf{Status final} \\
\midrule
1. Alucinação de citação & \textbf{CLEAR} — 0 inventadas; 17 refs conferidas \\
2. Bug de implementação & \textbf{CLEAR} — B$\geq$0, consistências, sensibilidade S23 ok \\
3. Resultados alucinados & \textbf{CLEAR} — todo número vem de script + CSV \\
4. Dependência de atalho & \textbf{CLEAR (sem flags)} \\
5. Bug-como-insight & \textbf{CLEAR} — achados contrários foram reportados, não racionalizados \\
6. Fabricação de metodologia & \textbf{CLEAR} — métodos padrão \\
7. Frame-lock de pipeline & \textbf{1 incidente} — frame ``fractal'' corrigido pelo autor, não auto-detectado \\
\bottomrule
\end{longtable}

% =====================================================================
\section{Avaliação da qualidade da colaboração}
\noindent\emph{Honestidade primeiro: sem inflar, com evidência. Avalia o desempenho do \textbf{autor}.}

\bigskip
\noindent\fbox{\parbox{0.96\linewidth}{%
\textbf{Pontuação de Qualidade da Colaboração: 82/100 — Excelente (faixa 75--89)}\\[6pt]
\begin{tabular}{@{}ll@{}}
Direção (Direction Setting)        & \scorebar{46.75}{85} \\[2pt]
Contribuição Intelectual           & \scorebar{48.4}{88} \\[2pt]
Controle de Qualidade (gatekeeping)& \scorebar{41.8}{76} \\[2pt]
Disciplina de Iteração             & \scorebar{46.75}{85} \\[2pt]
Eficiência de Delegação            & \scorebar{44}{80} \\[2pt]
Meta-Aprendizado                   & \scorebar{41.25}{75} \\
\end{tabular}}}

\bigskip
\noindent\textbf{Geral: 82/100.} O autor não foi um ``botão de continuar'': fez decisões direcionais
corretas e uma correção de rumo que elevou o trabalho.

\subsection*{O que funcionou bem}
\begin{itemize}\setlength\itemsep{1pt}
\item \textbf{O pivô} (\emph{``não é essa a pegada que eu quero fazer. Vamos voltar''}): reconhecer que um
  artigo tecnicamente bom não era o artigo certo, e redirecionar para a caracterização ancorada na
  tradição brasileira — a intervenção de maior valor. \textbf{(Contribuição Intelectual 88.)}
\item \textbf{Curadoria de literatura:} \emph{``estudasse a partir do Nereus\dots\ como o Haddad''} e
  \emph{``ancorar S3 em literatura nacional, o melhor do insumo''}.
\item \textbf{Disciplina:} re-rodar o pipeline (pivô), atacar os 4 obrigatórios \emph{e} os 4 sugeridos,
  investigar a anomalia de 2021. \textbf{(Iteração 85.)}
\item \textbf{Refinamento fino:} trocar ``fractal'' por termo da literatura; figuras em boas práticas.
\end{itemize}

\subsection*{Oportunidades perdidas}
\begin{itemize}\setlength\itemsep{1pt}
\item \textbf{Verificação independente dos números:} o autor confiou nos cálculos da IA sem checar
  amostras por conta própria — o gatekeeping numérico ficou a cargo dos gates da IA. \textbf{(Controle 76.)}
\item \textbf{Convenção de spillover/feedback} ficou pendente desde o projeto original; uma decisão
  explícita a fecharia.
\end{itemize}

\subsection*{Recomendações para a próxima vez}
\begin{enumerate}\setlength\itemsep{1pt}
\item Pedir à IA, num ponto, para listar os números-âncora e conferir 2--3 contra a fonte bruta.
\item Decidir cedo a convenção contestada (spillover/feedback).
\item Considerar \emph{cross-model} (\texttt{ARS\_CROSS\_MODEL}) nos gates para artigos de submissão.
\item Pedir o registro de lições nas skills (meta-aprendizado), não só no artigo.
\item Definir público-alvo/periódico no início — afina tom, extensão e exigência de benchmark.
\end{enumerate}

\subsection*{Valor humano vs.\ IA}
\begin{itemize}\setlength\itemsep{1pt}
\item \textbf{O frame.} A IA produziu um artigo ``fractal'' competente; \textbf{o autor escolheu o artigo
  certo} (caracterização). A direção intelectual foi humana.
\item \textbf{A âncora na tradição certa} (NEREUS/Haddad; desindustrialização nacional via IO) veio do
  conhecimento do autor.
\item A IA contribuiu com \textbf{velocidade, reprodutibilidade, rigor de verificação e crítica
  adversarial} — mas o \textbf{julgamento editorial} que define a identidade do artigo foi do autor.
\end{itemize}

\end{document}
